\begin{resumo}[Abstract]
    Dengue is one of the main public health challenges in Brazil, and the 2024 epidemic, with about 6.5 million notifications, highlighted the need for tools that anticipate outbreaks. This work developed and evaluated a model that forecasts weekly dengue notifications four weeks ahead for the 137 Brazilian mesoregions, integrating SINAN microdata with data from about 700 INMET automatic weather stations through a reproducible layered pipeline. A data audit identified and corrected defects in the public files, including a shift in the column mapping of the INMET files that made the variable treated as humidity correspond to the dew point, and showed that part of the preliminary results were artifacts of these defects. Three XGBoost variants, with and without climate variables, were compared with naive forecasters on a 2024--2026 test set. The model with climate reached an $R^2$ of 0.87 on the logarithmic scale, above all naive forecasters, and outperformed persistence on every metric in 2025 and 2026; in the 2024 epidemic, however, it anticipated the timing of the peak but predicted about half of the notifications, and persistence was better on the original scale. Climate variables provided a small but statistically consistent gain, present in every temporal validation window and in scenarios without recent epidemiological data, and on their own explained 39\% of the variance on the logarithmic scale. For risk-level classification, persistence outperformed dedicated classifiers. The results were made available through a REST API and visualization dashboards.

     \vspace{\onelineskip}

     \noindent
     \textbf{Keywords}: Dengue. Forecasting. Machine learning. XGBoost. Climate. Epidemiological surveillance.
\end{resumo}
